% ============================================================
% DRAFT ONLY — not meant to compile, nothing \input's this file.
%
% WHAT THIS IS: edit-by-edit instructions for tightening
% design.tex, covering the subsections you flagged (State Machine,
% Failure Recovery, Conditional Commands) plus the Table 3.1 trim
% we discussed. I left Plan Check, Confirmation and Later Moments
% untouched since you said those are fine.
%
% Format per edit: CURRENT (commented, quoted from the live file so
% you can find it with Ctrl+F) then REPLACE WITH (not commented —
% paste that text in directly). Section headers and \label{} lines
% are never touched, only body text.
%
% ONE EDIT NEEDS A COMPANION CHANGE IN implementation.tex — marked
% clearly below (EDIT 3). Skipping that half would genuinely lose a
% technical detail (the halt() flag mechanism), not just shorten
% something, so don't skip it.
% ============================================================


% ============================================================
% EDIT 1 — Table 3.1 (User and functional requirements)
% Drop the "Objective" column (it only re-points at the Objectives
% list, adds nothing new) and tighten the requirement wording. Keeps
% the FR/NR IDs intact since NR1, NR3 and FR9 are cited elsewhere in
% this chapter (System Architecture and Reporting) — cutting the
% table entirely would orphan those three citations.
%
% CURRENT: the whole \begin{table}[htbp] ... \end{table} environment
% under \label{tab:requirements} (the one with the "ID / Requirement
% / Objective" header row).
%
% REPLACE WITH:

\begin{table}[htbp]
  \centering
  \caption{Functional (FR) and non-functional (NR) requirements.}
  \label{tab:requirements}
  \small
  \begin{tabular}{@{}l p{10.8cm}@{}}
    \toprule
    ID & Requirement \\
    \midrule
    FR1 & Take the command by voice and transcribe it locally. \\
    FR2 & Check the command against the room before planning, and ask a spoken question when it is ambiguous, up to a fixed number of questions. \\
    FR3 & Produce a plan that uses only actions the rover has and objects the room contains. \\
    FR4 & Read the plan back and move only after a spoken confirmation; a rejection returns to the dialogue rather than ending it. \\
    FR5 & Stop without consulting a model when an obstacle blocks a step. \\
    FR6 & Use the camera to check progress and arrival during a step. \\
    FR7 & Report whether the rover completed the steps and flag anything it was unsure about. \\
    FR8 & On a failure (never on ambiguity), ask the person what to do next and carry on from where the rover is. \\
    FR9 & Log every model call and decision made, to support debugging. \\
    FR10 & On a conditional command, decide the condition once, run only the matching steps, and ask the person when the camera cannot tell. Never guess a branch. \\
    FR11 & When asked what is near a target, report what surrounds it, not only whether the target was reached. \\
    NR1 & Run planning, vision and speech on consumer-grade hardware (e.g.\ a Ryzen 5 3600, 16\,GB RAM, 12\,GB GPU). Cloud models are allowed only as a benchmark. \\
    NR2 & Keep the rover stoppable when the network or a model is slow or unavailable. \\
    NR3 & Use only a camera, an ultrasonic sensor and a line sensor — no encoders, IMU or lidar. \\
    NR4 & Let the planner or vision model be swapped by configuration, with no change to the pipeline code. \\
    NR5 & Show the rover's state without a screen, using onboard lights. \\
    \bottomrule
  \end{tabular}
\end{table}


% ============================================================
% EDIT 2 — State Machine (\label{ssec:statemachine})
% Cuts one sentence that just restates "capped at five questions"
% a second time, and fixes typos along the way (wether/whether,
% sued/used, quetion/question, coould/could, textt/texttt,
% "I am no fully sure I understoon"/"I am not fully sure I
% understood").
%
% CURRENT: the full body of \mysubsection{The dialogue as a state
% machine}{State Machine}, from "The command spoke moves through..."
% down to "...could result in letting every command through."
%
% REPLACE WITH:

The command moves through the phases in Figure~\ref{fig:states}. In \texttt{CLARIFYING}, the
vision model is asked whether the command is ambiguous, given the command and the scene
digest, and answers in JSON with four fields: whether it is ambiguous, why, a question to ask,
and a rewritten command. An ambiguous command gets the question spoken and the answer added to
the history; a clear command uses the rewritten version from then on, which only names objects
and avoids words like `left' or `behind you', since those are relative to the camera's
position when asked and can be wrong once the rover moves.

The loop is capped at five questions, after which the session is marked capped and planning
proceeds from whichever command is available — rewritten if a question was answered, original
if none was. A question already asked is dropped from the next round; if that leaves nothing
new to ask, the loop breaks early onto planning. The read-back plan is prefixed with a warning
(`I am not fully sure I understood') whenever the session was capped, so the person knows the
plan is not certain.

If the vision model fails to load, it is retried once; a second failure tells the person and
asks them to repeat the command, rather than letting it through unchecked.


% ============================================================
% EDIT 3 — Failure Recovery (\label{sec:recovery})
% THREE separate changes in this section. The third one MOVES a
% paragraph to implementation.tex rather than cutting it — see the
% companion edit at the bottom of this block. Do all three or the
% halt() mechanism disappears from the report entirely.

% --- 3a: opening paragraph ---
% CURRENT: "From the pre-video recording which generates the scene
% digest. New objects can be introduced into the room or old objects
% can be removed. Situations like an 'orange box' being behind a
% chair or a bag is left in the rovers path. If the rover reached the
% desired destination the plan named but the camera sees a different
% scene. This is a case where the plan was correct but the world did
% not coporate."
%
% REPLACE WITH:

The pre-recorded scene digest can go stale: an object can be moved, removed, or something new
(a bag left in the rover's path) can appear by the time the rover gets there. If the rover
reaches the location the plan named but the camera sees a different scene, the plan was
correct but the world did not cooperate.

% --- 3b: the guidance-response paragraph ---
% CURRENT: "The person can then repond by answering what the rover
% should do next.  The planner can then write a new plan for what
% it still has to do given  the original command, the steps that it
% already completed, the steps that failed and why it failed. What
% the rover has already seen and the conversation up until that
% point. The prompt then tells the planner to use the answer to
% decide where to look or what to reach instead. If the person says
% where the target is, the rover should turn that way and approach
% the target again."
%
% REPLACE WITH:

The person responds with what the rover should do next. The planner then writes a new plan
using the original command, the steps already completed, the failed step and why, what the
rover has seen, and the conversation so far, with the person's answer used to decide where to
look or what to approach instead — if they name where the target is, the rover turns that way
and tries the approach again.

% --- 3c: the halt() / TurboPi flag paragraph — MOVE, don't cut ---
% CURRENT (last paragraph of the section): "The \texttt{halt()}
% function is not called in a recovery stage. If a step is blocked
% it stops counting the time it tales to execute that step.The
% turboPi then sets a \texttt{halt()} flag that only the next
% \texttt{execute\_step} clears. Since the free-roam approach in
% Section~\ref{ssec:freeroam} directly calls the drive commands on
% the rover, a flag that is not cleared could cause issues down the
% line."
%
% This is implementation detail, not design reasoning — it belongs
% in implementation.tex next to the other recovery-loop code detail,
% not here. REPLACE the paragraph above WITH this one line in
% design.tex:

Entering recovery does not itself call \texttt{halt()} — how the rover is kept stopped across
a blocked step is implementation detail, covered in Section~\ref{ssec:guidance_loop}.

% AND separately, add this new paragraph to implementation.tex, at
% the end of \mysubsection{Failure reovery in code}{Failure Recovery
% in Code} (\label{ssec:guidance_loop}) — it isn't there yet:

\texttt{halt()} is not called when entering recovery itself. A blocked step stops the rover's
step-timing clock, and the TurboPi sets a \texttt{halt()} flag that only the next
\texttt{execute\_step} call clears. This matters because the free-roam approach
(Section~\ref{ssec:freeroam}) calls the drive commands directly, so a flag left set could
otherwise block the next approach step from moving at all.


% ============================================================
% EDIT 4 — Conditional Commands (\label{sec:conditional})
% Four small tightenings through the section; typos fixed along the
% way (decisioins/decisions, seperated/separated, repond/respond,
% continueally/continually, "Rather than guaranteeing" fragment-style
% phrasing cleaned up where it occurs here too).

% --- 4a ---
% CURRENT: "A conditional command, which is level L4 test in
% Table~\ref{tab:command-levels},asks the rover to make a decision on
% something and act on the answer: 'check whether the orange box is
% visible, and if it is, go to it, otherwise stop'. Who actually
% makes the decision in this case. The planner writes both branches
% up front, before anything moves, and the rover only evaluates the
% one condition at the moment it matters. This allows the planner to
% do the thing it is good at, and turning words into structured plan.
% Then leaving the single judgements that need a vision model call to
% the camera and vision model."
%
% REPLACE WITH:

A conditional command (level L4 in Table~\ref{tab:command-levels}) asks the rover to decide
something and act on the answer, e.g.\ `check whether the orange box is visible, and if it is,
go to it, otherwise stop'. The planner writes both branches up front, before anything moves,
and the rover evaluates only the one condition that matters at the time — letting the planner
do what it is good at, turning words into a structured plan, while leaving the single judgement
that needs a camera and vision model to the point where it actually matters.

% --- 4b ---
% CURRENT: "When the rover reaches a \texttt{check} step in the plan.
% The rover centers the gimbal, waits half a second for it to settle,
% and then takes one frame (Figure~\ref{fig:cond}). The vision model
% is asked whether the stated condition is true of what is visible
% and it can answer yes, no or unsure. Now this loop can remembers
% the answer. Every later step tagged with that answer is skipped and
% logged as skipped. The run record now can show which branch ran and
% which did not."
%
% REPLACE WITH:

At a \texttt{check} step, the rover centres the gimbal, waits half a second to settle, and
takes one frame (Figure~\ref{fig:cond}). The vision model answers yes, no or unsure for the
stated condition; every later step tagged with the opposite answer is skipped and logged as
skipped, so the run record shows which branch ran.

% --- 4c ---
% CURRENT: "Now the choice of what happens on the 'unsure', and on a
% failed vision call. The step returns \texttt{blocked}, and the
% existing recovery loop (Section~\ref{sec:recovery}) asks the person
% what to do. There is no default branch added and if the rover
% guesses 'no' continually it would prompt the rover to go to the
% incorrect object continueally without anyone understanding if it
% was working."
%
% REPLACE WITH:

An `unsure' answer, or a failed vision call, returns \texttt{blocked} and goes to the existing
recovery loop (Section~\ref{sec:recovery}) rather than defaulting to a branch — a rover that
silently guessed `no' every time could run for a long time without anyone noticing it wasn't
working.

% --- 4d ---
% CURRENT: "There are limits, and I would state them here. The check
% is one frame from wherever the rover is standing. It does not
% search. A condition about an object just outside the camera view
% will come back with a 'no' or 'unsure' even if the object is in the
% room. The results depends on where the rover is facing."
%
% REPLACE WITH:

The check is limited to one frame from wherever the rover is standing — it does not search, so
a condition about an object just outside the camera's view can come back `no' or `unsure' even
though the object is in the room; the result depends on which way the rover happens to be
facing.

% ============================================================
% Rough page impact: Edit 1 ~0.3 page, Edit 2 ~0.2 page, Edit 3
% ~0.15 page (plus the moved paragraph, net page count for Design
% unaffected by the move itself since it's leaving the chapter, not
% getting cut — implementation gains it back), Edit 4 ~0.2 page.
% Call it under a page total from these four. Getting from ~18 pages
% to 10 still needs the flow-diagrams-to-appendix move and a similar
% pass over System Architecture and Free-roam Approach, which we
% haven't touched yet — flag if you want those done the same way.
% ============================================================
